\chapter{Going 3-D: the same code, a different economy}
\label{ch:A5_three_dimensions}
\noindent\texttt{tutorials/A\_foundations/A5\_three\_dimensions.py}\\[0.75em]
\begin{outcome}
You can run every A-track idea in three dimensions by
changing ONE argument (dim=3) — and you can articulate what actually
changed: the cost model. Elements per level grow 8x (not 4x), DOFs per p2
element go 9 -> 27, and the direct solver's fill-in makes the SOLVE the
bottleneck far earlier. Dimension-independence of the CODE is a design
property; dimension-independence of the COST is a myth.
\end{outcome}

The whole stack is written over k-D octrees (this library also
runs k=4 for future space-time work; the test suite exercises it). This
chapter: (i) the A1 MMS in 3-D, orders 2 and 3 again; (ii) the A3 immersed
solve on a SPHERE, order 2 again; (iii) a small cost table you must extend.

\begin{expected}
\begin{verbatim}
box  p1: orders 2.00 / 2.00      box  p2: orders 2.97 / 2.99
    sphere-immersed p1: order 1.34 at the (3,4) level pair — PREASYMPTOTIC:
    at level 3 the sphere is ~5 elements across and the surrogate staircase
    is still coarse geometry, not discretization. The resolved pair (4,5)
    gives ~2.0 (the test suite locks it); running it here costs minutes —
    that is Explore (c)'s point about 3-D economics.
    (levels are one lower than 2-D — 8x growth per level.)
\end{verbatim}
\end{expected}

\section*{The script}
\begin{lstlisting}
import time

import numpy as np
from scipy.sparse.linalg import splu
import warp as wp

from diffsim.octree.build import build_uniform
from diffsim.mesh.nodes import build_mesh
from diffsim.mesh.constraints import build_constraints
from diffsim.mesh.basis import basis_tables
from diffsim.mesh.faces import face_tables
from diffsim.assembly.operators import DeviceMesh, assemble_csr
from diffsim.geometry.csg import Sphere
from diffsim.sbm.surrogate import classify_lambda, extract_surrogate, GeometryData
from diffsim.sbm.poisson import SBMPoisson
from diffsim.physics.poisson import (gauss_points, make_load_kernel,
                                     l2_error_masked)

DEVICE = "cuda:0"
u3 = lambda x: (np.sin(np.pi * x[:, 0]) * np.sin(np.pi * x[:, 1])
                * np.sin(np.pi * x[:, 2]))
f3 = lambda x: 3 * np.pi ** 2 * u3(x)


def solve_box(level, p):
    tree = build_uniform(level, dim=3)
    mesh = build_mesh(tree, p=p)
    cons = build_constraints(mesh)
    dm = DeviceMesh.from_mesh(mesh, cons, basis_tables(p, dim=3), DEVICE)
    A = assemble_csr(dm)
    F_full = wp.zeros(dm.n_nodes, dtype=wp.float64, device=DEVICE)
    xq = gauss_points(mesh, dm.tables_by_p)
    for pv, b in dm.bins.items():
        fq = wp.array(f3(xq[pv]), dtype=wp.float64, device=DEVICE)
        lk = make_load_kernel(b["nbf"], b["nqp"], dm.dim)
        wp.launch(lk, dim=len(b["eids"]),
                  inputs=[b["conn"], b["h"], b["N"], b["w"], fq, F_full],
                  device=DEVICE)
    b_vec = np.asarray(cons.T.T @ F_full.numpy())
    coords = mesh.node_coords[cons.free_nodes]
    bdry = np.where(mesh.boundary_nodes[cons.free_nodes])[0]
    A = A.tolil()
    for i in bdry:
        A.rows[i] = [int(i)]
        A.data[i] = [1.0]
        b_vec[i] = 0.0                       # u* = 0 on the box boundary
    u_free = splu(A.tocsr().tocsc()).solve(b_vec)
    u_all = np.asarray(cons.T @ u_free)
    return l2_error_masked(dm, u_all, u3, lambda x: np.ones(len(x), bool))


def solve_sphere(level):
    oracle = Sphere((0.5, 0.5, 0.5), 0.3)
    tree = build_uniform(level, dim=3)
    ret, _ = classify_lambda(tree, oracle, 0.0)
    sf = extract_surrogate(ret)
    mesh = build_mesh(ret, p=1)
    cons = build_constraints(mesh)
    dm = DeviceMesh.from_mesh(mesh, cons, basis_tables(1, dim=3), DEVICE)
    geo = GeometryData.evaluate(oracle, ret, sf, face_tables(1, 3))
    u = SBMPoisson(dm, geo, sf, g_fn=u3).solve(f_fn=f3)
    return l2_error_masked(dm, u, u3, lambda x: oracle.classify(x) < 0)


if __name__ == "__main__":
    for p, levels in ((1, (3, 4, 5)), (2, (2, 3, 4))):
        errs = [solve_box(lv, p) for lv in levels]
        orders = [np.log2(errs[i] / errs[i + 1]) for i in range(2)]
        print(f"box  p={p}: errors " + "  ".join(f"{e:.3e}" for e in errs)
              + "   orders " + "  ".join(f"{o:.2f}" for o in orders))
    errs = [solve_sphere(lv) for lv in (3, 4)]
    print(f"sphere p=1: errors " + "  ".join(f"{e:.3e}" for e in errs)
          + f"   order {np.log2(errs[0] / errs[1]):.2f}")
    # a first cost row for the Explore table
    t0 = time.time(); solve_box(4, 1); t1 = time.time()
    print(f"\nbox level 4 p1 total: {t1 - t0:.2f} s  <- extend this table")
\end{lstlisting}

\begin{explore}
\begin{verbatim}
(a) Extend the cost table: levels 3..6, stages separated (mesh+
      constraints / assemble / solve). In 2-D the splu solve scales
      ~O(n^1.5); in 3-D it is ~O(n^2) with O(n^{4/3}) memory — find the
      level where the solve overtakes assembly, and compare with the 2-D
      crossover from A1's performance corner. This is THE reason track-P's
      iterative solvers exist.
  (b) The suite runs the same physics at k=4 (see tests/, 'dim-coverage').
      Estimate — do not run — the level-4 k=4 element count and p2 DOF
      count. At what level would k=4 fit in your GPU memory?
  (c) Rerun the sphere at level 5. Time it. Where did the time go —
      classification, constraints, assembly, or solve? (Profile before
      you answer; the m0.5 findings say constraints, and explain why.)
\end{verbatim}
\end{explore}
